\section{Rigorous Derivations of Blocking Lemmas}
\label{app:blocking_lemmas}

Based on the document, to complete the proof of the Yang-Mills Mass Gap, we provide the rigorous mathematical derivations for four specific "blocking lemmas." These are the technical bridges that turn the physical framework into a mathematical theorem.

\subsection{Derivation: The Lee-Yang Zero Bound (Absence of Phase Transitions)}

\textbf{Goal:} Prove that the mass gap $\Delta > 0$ does not vanish as you interpolate from Supersymmetric Yang-Mills ($N_f=1$) to Pure Yang-Mills ($N_f=0$).

\textbf{Requirement:} You must derive a zero-free region for the partition function in the complex mass plane to ensure analyticity.

\subsubsection*{Mathematical Steps}

\begin{enumerate}
    \item \textbf{Define the Partition Function:}
    Let $Z(m) = \int [dA] \det(D + m) e^{-S_{YM}}$.
    For Adjoint QCD, the Dirac operator $D$ is anti-Hermitian. Its eigenvalues are purely imaginary $i\lambda_k$ with $\lambda_k \in \mathbb{R}$.

    \item \textbf{Factorize the Determinant:}
    Show that the determinant is real and positive for real $m$:
    \begin{equation}
        \det(D+m) = \prod_k (i\lambda_k + m) = \prod_{k>0} (m^2 + \lambda_k^2) > 0
    \end{equation}

    \item \textbf{Establish the Zero-Free Strip:}
    Derive the bound on the location of zeros $m_0$. Since zeros occur only at $m = \pm i \lambda_k$, and the lattice Dirac operator is bounded by $\|D\| \le \frac{c}{a}$ (where $a$ is lattice spacing), all zeros lie on the imaginary axis segment $[-i\frac{c}{a}, i\frac{c}{a}]$.

    \item \textbf{Apply Complex Pirogov-Sinai Theory:}
    To rule out zeros pinching the real axis in the infinite volume limit (which would signal a phase transition), construct "contours" (regions where the field deviates from the vacuum).
    \begin{itemize}
        \item Define the contour activity $\mathfrak{z}(\Gamma)$.
        \item \textbf{Derive the Peierls Condition:} Prove that for complex mass $m$ with $\text{Re}(m) > m_0$, the real part of the string tension provides exponential decay:
        \begin{equation}
            |\mathfrak{z}(\Gamma)| \le e^{-\sigma \text{Area}(\Gamma)}
        \end{equation}
        \item \textbf{Convergence:} Use this bound to satisfy the Kotecký-Preiss criterion:
        \begin{equation}
            \sum_{\Gamma \nsim \Gamma_0} |\mathfrak{z}(\Gamma)| e^{a |\Gamma|} \le 1
        \end{equation}
        proving that the free energy is analytic in the strip $|\text{Im}(m)| < \delta$.
    \end{itemize}
\end{enumerate}

\subsection{Derivation: The Battle-Federbush Bound (UV Stability)}

\textbf{Goal:} Prove that adding fermions does not destroy the ultraviolet stability established by Balaban for pure gauge theory.

\textbf{Requirement:} Derive a bound on the fermion determinant that allows it to be controlled by the gauge action.

\subsubsection*{Mathematical Steps}

\begin{enumerate}
    \item \textbf{Define the Regularized Determinant:}
    The standard determinant diverges in 4D. Use the Hilbert-Carleman regularized determinant $\det_p(1+K)$ where $K = D^{-1} \delta A$:
    \begin{equation}
        \ln \det{}_p (1+K) = \text{Tr} \left( \ln(1+K) - \sum_{j=1}^{p-1} \frac{(-1)^{j-1}}{j} K^j \right)
    \end{equation}
    For 4D, choose $p=5$ to subtract diagrams with divergence degree $\ge 0$.

    \item \textbf{Cluster Expansion on the Determinant:}
    Expand the determinant into a sum over tree graphs connecting local operator insertions (Battle-Federbush expansion).

    \item \textbf{Derive the Tree Decay Bound:}
    Prove that the kernel $K(x,y)$ decays exponentially with distance $|x-y|$ (using the Combes-Thomas estimate).
    \begin{equation}
        |K(x,y)| \le C e^{-m |x-y|}
    \end{equation}

    \item \textbf{Compare to Gauge Action:}
    Show that the logarithm of the determinant is bounded by the kinetic energy of the gauge field:
    \begin{equation}
        |\ln \det{}_p (1+K)| \le \epsilon \|F_{\mu\nu}\|^2
    \end{equation}
    This ensures that the total effective action $S_{eff} = S_{YM} - \ln \det$ remains bounded from below (stable).
\end{enumerate}

\subsection{Derivation: The 1D Transfer Matrix Gap (Uniform LSI Base Case)}

\textbf{Goal:} Establish the base case for the Hierarchical Log-Sobolev Inequality.

\textbf{Requirement:} Explicitly compute the spectral gap for a 1D chain of $SU(N)$ matrices.

\subsubsection*{Mathematical Steps}

\begin{enumerate}
    \item \textbf{Define the Transfer Operator:}
    \begin{equation}
        (T \psi)(U) = \int dV e^{\beta \text{Re Tr}(U V^\dagger)} \psi(V)
    \end{equation}

    \item \textbf{Peter-Weyl Decomposition:}
    Expand the kernel in characters $\chi_r(U)$. The operator is diagonal in the representation basis. The eigenvalues are ratios of character expansion coefficients:
    \begin{equation}
        \lambda_r(\beta) = \frac{I_{d_r}(\beta)}{I_0(\beta)}
    \end{equation}

    \item \textbf{Identify the Gap:}
    The spectral gap is determined by the fundamental representation $F$:
    \begin{equation}
        \text{Gap} = 1 - \lambda_F = 1 - \frac{I_1(\beta)}{I_0(\beta)}
    \end{equation}

    \item \textbf{Derive the Turán Inequality:}
    For $SU(2)$, the coefficients are modified Bessel functions $I_\nu(\beta)$. You must derive the bound:
    \begin{equation}
        \frac{I_2(\beta)}{I_1(\beta)} < \frac{I_1(\beta)}{I_0(\beta)} \implies \lambda_{F} < 1
    \end{equation}

    \item \textbf{Calculate the Explicit Lower Bound:}
    Use asymptotic expansions of $I_\nu(\beta)$ to derive the uniform bound required for the induction:
    \begin{equation}
        \text{Gap} \ge \frac{1}{\beta} - O(\frac{1}{\beta^2})
    \end{equation}
\end{enumerate}

\subsection{Derivation: The Giles-Teper Bound (Spectral Lower Bound)}

\textbf{Goal:} Convert the string tension $\sigma$ (area law) into a mass gap $M$ (eigenvalue gap).

\textbf{Requirement:} Derive the inequality $M \ge \sqrt{\sigma}$ using the variational principle.

\subsubsection*{Mathematical Steps}

\begin{enumerate}
    \item \textbf{Variational Characterization:}
    Recall the mass gap definition $M = \inf_{\Psi \perp \Omega} \frac{\langle \Psi, H \Psi \rangle}{\langle \Psi, \Psi \rangle}$.

    \item \textbf{Construct the Flux Tube State:}
    Define a trial state $\Psi_L$ corresponding to a flux loop of length $L$.

    \item \textbf{Evaluate Energy Expectation:}
    Using the Transfer Matrix spectral decomposition, show the energy of this state is:
    \begin{equation}
        E(L) \approx \sigma L + \frac{c}{L}
    \end{equation}

    \item \textbf{Optimize the Geometry:}
    Minimize the energy density with respect to the loop size $L$. The kinetic term (from localization of the flux) scales as $1/L$.
    \begin{equation}
        \frac{d}{dL} \left( \sigma L + \frac{1}{L} \right) = 0 \implies L_{opt} = \frac{1}{\sqrt{\sigma}}
    \end{equation}

    \item \textbf{Final Inequality:}
    Substitute $L_{opt}$ back into the energy equation to derive the bound:
    \begin{equation}
        M \ge 2\sqrt{\sigma}
    \end{equation}
    \textit{Note: To get the \textbf{lower} bound $M \ge \sqrt{\sigma}$, you must use the \textbf{Cheeger Inequality} on the gauge orbit space configuration manifold, relating the geometric bottleneck (string tension) to the Laplacian spectral gap.}
\end{enumerate}
